\par\noindent\begin{minipage}{\linewidth}
% !TeX root = ../教材题库_全部条目_转写工作稿.tex
% Source: XJ-B1-C01-SU-P028-Y1490
% 第1章 / 小结与复习 小结与复习 / 教材 p.28 / PDF 第35页
% 已对照原页转写题干及公式；待教师复核；未制作答案
% 原题及相邻原解截图（核对用）：../题目原图/XJ-B1-C01-SU-P028-Y1490-01.webp
% 本文件只在转写完成后作为题干使用；OCR 条目尚不是可靠题干。
\begin{bankitem}{第1章 小结与复习 第6题}{第1章 / 小结与复习 小结与复习 / 教材 p.28 / PDF 第35页}{XJ-B1-C01-SU-P028-Y1490}{已对照原页转写题干及公式；待教师复核；未制作答案}
当陈述语句中涉及不确定的对象时，其真假往往难以确定。使用全称量词和存在量词可以对不确定的对象进行约束，从而使语句具有确定的真假。试结合实例讨论：如何判断含有量词的命题的真假？如何对含有一个量词的命题进行否定？
\end{bankitem}
\end{minipage}\par\medskip
